\chapter{分子转子、不可约张量与激光排列}
\label{ch:molecular-rotation}

双原子分子的轴只需两个角定位，一般刚性多原子分子还必须说明绕自身轴转了多少，因此构型空间是整个 \(SO(3)\)。三个主转动惯量把刚体分成球陀螺、对称陀螺和非对称陀螺；对称陀螺保留体轴投影 \(K\)，非对称性则把相差二的 \(K\) 分量混合。空间固定与体固定角动量分别来自旋转群的左右作用，它们看似相反的对易号正是左右作用的几何结果。分子极化率再作为秩二球张量与激光相连，选择定则和脉冲后的转动回生便从同一表示结构中产生。
\section{SO(3) 上的量子转子与 Wigner 基}

刚体的“位置”不是质点坐标，而是它相对于实验室坐标系的取向。选定一个参考取向以后，任意新取向由唯一的三维适当旋转给出，因此构型空间是
\[
SO(3)=\{R\in M_3(\mathbb R):R^{\mathsf T}R=I,\ \det R=1\}.
\]
矩阵 \(R\) 的三列是分子三根体轴在实验室坐标中的分量。条件
\(R^{\mathsf T}R=I\) 保证三根轴仍正交归一，\(\det R=1\) 保证取向没有被镜面反射。

\begin{definition}[群上的波函数与左右作用]
量子转子的波函数是平方可积函数
\(\psi:SO(3)\to\mathbb C\)。对 \(S\in SO(3)\)，定义
\[
(L_S\psi)(R)=\psi(S^{-1}R),\qquad
(R_S\psi)(R)=\psi(RS).
\]
\(L_S\) 改变实验室参考系中的取向，\(R_S\) 改变分子本体参考系中的取向。两种作用满足
\(L_{S_1}R_{S_2}=R_{S_2}L_{S_1}\)。
\end{definition}

最后一个等式直接来自矩阵乘法结合律：
\[
(L_{S_1}R_{S_2}\psi)(R)
=\psi(S_1^{-1}RS_2)
=(R_{S_2}L_{S_1}\psi)(R).
\]
因此空间固定角动量与体固定角动量可以同时对角化。

\begin{definition}[Haar 测度]
若群上的积分测度 \(\dd R\) 满足
\[
\int f(SR)\,\dd R=\int f(R)\,\dd R
=\int f(RS)\,\dd R
\]
对所有 \(S\in SO(3)\) 成立，就称为 Haar 测度。紧群上的 Haar 测度在整体常数因子之外唯一。采用 Euler 角
\(R=R_z(\alpha)R_y(\beta)R_z(\gamma)\) 时，可取
\[
\dd R=\dd\alpha\,\sin\beta\,\dd\beta\,\dd\gamma,
\qquad
\int_{SO(3)}\dd R=8\pi^2.
\]
\end{definition}

因子 \(\sin\beta\) 不是人为插入的权重。它与球坐标面积元中的
\(\sin\theta\) 来源相同：在 Euler 角坐标变换下，群流形的体积元必须乘以相应 Jacobian。左乘或右乘只重新参数化同一群流形，所以这个体积元保持不变。

\begin{definition}[Wigner 矩阵元]
\(SO(3)\) 的角动量 \(J\) 不可约表示记为 \(D^J(R)\)。在
\(J_z\) 本征基中，它的矩阵元
\[
D^J_{MK}(R)=\langle JM|D^J(R)|JK\rangle
\]
称为 Wigner \(D\) 函数。指标 \(M\) 属于空间固定基，指标
\(K\) 属于体固定基。
\end{definition}

在分子主惯量轴 \(a,b,c\) 中，惯量张量已经对角化：

\[
\mathbf I
=\operatorname{diag}(I_a,I_b,I_c).
\]

经典角速度在体固定坐标中的分量记为 \(\boldsymbol\Omega=(\Omega_a,\Omega_b,\Omega_c)\)。转动动能和角动量分别为

\[
T_{\rm rot}
=\frac12\boldsymbol\Omega^T\mathbf I\boldsymbol\Omega,
\qquad
J_i=I_i\Omega_i.
\]

消去 \(\Omega_i\) 后得到

\[
T_{\rm rot}
=\frac{J_a^2}{2I_a}
+\frac{J_b^2}{2I_b}
+\frac{J_c^2}{2I_c}.
\]

正则量子化后，刚性多原子转子的哈密顿量为

\begin{equation}
H_{\rm rot}
=\frac{\hat J_a^2}{2I_a}
+\frac{\hat J_b^2}{2I_b}
+\frac{\hat J_c^2}{2I_c}.
\label{eq:asymmetric-top-hamiltonian}
\end{equation}

定义能量量纲的主轴转动常数

\[
A_{\rm rot}=\frac{\hbar^2}{2I_a},
\qquad
B_{\rm rot}=\frac{\hbar^2}{2I_b},
\qquad
C_{\rm rot}=\frac{\hbar^2}{2I_c},
\]

并用无量纲角动量 \(\mathbf j=\mathbf J/\hbar\)，则

\begin{equation}
H_{\rm rot}
=A_{\rm rot}j_a^2+B_{\rm rot}j_b^2+C_{\rm rot}j_c^2.
\label{eq:asymmetric-top-energy-form}
\end{equation}

球形陀螺满足 \(I_a=I_b=I_c\)；长球对称陀螺满足 \(I_a<I_b=I_c\)，即 \(A_{\rm rot}>B_{\rm rot}=C_{\rm rot}\)；扁球对称陀螺满足 \(I_a=I_b<I_c\)，即 \(A_{\rm rot}=B_{\rm rot}>C_{\rm rot}\)。三者都不相等时才是真正的不对称陀螺。

刚体取向由三个 Euler 角 \((\alpha,\beta,\gamma)\) 描述。旋转群 \(SO(3)\) 上的一组自然正交基是 Wigner 矩阵元

\[
D^J_{MK}(\alpha,\beta,\gamma),
\]

其中 \(J\) 是总角动量量子数，\(M\) 是角动量在空间固定 \(Z\) 轴上的投影，\(K\) 是在体固定对称轴上的投影。规范化转动态可记为

\[
\langle\alpha,\beta,\gamma|JKM\rangle
=\sqrt{\frac{2J+1}{8\pi^2}}\,D^{J*}_{MK}(\alpha,\beta,\gamma),
\]

其中归一化采用 Euler 角测度 \(\dd\alpha\,\sin\beta\,\dd\beta\,\dd\gamma\)，且 \(0\le\alpha,\gamma<2\pi\)、\(0\le\beta\le\pi\)。

并满足

\begin{align*}
\mathbf j^2|JKM\rangle&=J(J+1)|JKM\rangle,\\
j_Z|JKM\rangle&=M|JKM\rangle,\\
j_a|JKM\rangle&=K|JKM\rangle
\end{align*}

（最后一式以长球对称轴 \(a\) 为例）。

对于长球对称陀螺，\eqnrefs{eq:asymmetric-top-energy-form} 化为

\[
H
=B_{\rm rot}\mathbf j^2
+(A_{\rm rot}-B_{\rm rot})j_a^2,
\]

所以 \(|JKM\rangle\) 已经是能量本征态，

\begin{equation}
E_{JK}^{\rm pro}
=B_{\rm rot}J(J+1)
+(A_{\rm rot}-B_{\rm rot})K^2.
\label{eq:prolate-symmetric-top-spectrum}
\end{equation}

类似地，扁球对称陀螺取对称轴 \(c\)，得到

\begin{equation}
E_{JK}^{\rm obl}
=B_{\rm rot}J(J+1)
+(C_{\rm rot}-B_{\rm rot})K^2.
\label{eq:oblate-symmetric-top-spectrum}
\end{equation}

无外场时哈密顿量不依赖空间固定投影 \(M\)，因此保持 \(2J+1\) 重简并；对于 \(K\neq0\)，能量又只依赖 \(K^2\)，于是 \(\pm K\) 成对简并。

\begin{proof}[解]
令 \(D^J(R)\) 是 \(SO(3)\) 的 \(2J+1\) 维不可约酉表示，并采用
\[
 \dd R=\dd\alpha\,\sin\beta\,\dd\beta\,\dd\gamma,
 \qquad \int_{SO(3)}\dd R=8\pi^2 .
\]
对固定的 \(K,K'\) 构造从 \(J'\) 表示空间到 \(J\) 表示空间的算符
\begin{equation*}
 X
 =\int\dd R\,
 D^J(R)\,|K\rangle\langle K'|\,D^{J'}(R)^\dagger.
\end{equation*}
Haar 测度的左不变性给出
\begin{equation*}
 D^J(R_0)X
 =\int\dd R\,D^J(R_0R)|K\rangle\langle K'|D^{J'}(R)^\dagger
 =XD^{J'}(R_0),
\end{equation*}
所以 \(X\) 是交织算符。Schur 引理要求 \(J\ne J'\) 时
\(X=0\)，而 \(J=J'\) 时 \(X=cI\)。取迹可得
\[
 (2J+1)c
 =\int\dd R\,\operatorname{tr}(|K\rangle\langle K'|)
 =8\pi^2\delta_{KK'} .
\]
最后取 \(M,M'\) 矩阵元，得到
\begin{equation}\label{eq:wigner-d-haar-orthogonality}
 \int\dd R\,
 D^{J*}_{MK}(R)\,D^{J'}_{M'K'}(R)
 =\frac{8\pi^2}{2J+1}
 \delta_{JJ'}\delta_{MM'}\delta_{KK'},
\end{equation}
\end{proof}

因此
\(\sqrt{(2J+1)/(8\pi^2)}\,D^{J*}_{MK}\)
是一组正交归一基；Peter--Weyl 定理保证它在
\(L^2(SO(3))\) 中完备。左、右正则作用彼此对易，使实验室系投影
\(M\) 与分子本体系投影 \(K\) 可以同时标记基矢。对称陀螺的
哈密顿量还与相应的体固定生成元对易，所以 \(K\) 守恒；
不对称项则只保留 \(K\) 的奇偶性。

\section{左右生成元与不对称陀螺}
选择体固定 \(c\) 轴作为量子化轴，定义无量纲升降算符

\[
j_\pm=j_a\pm\ii j_b.
\]

利用

\[
j_a=\frac{j_++j_-}{2},
\qquad
j_b=\frac{j_+-j_-}{2\ii},
\]

有

\begin{align*}
A_{\rm rot}j_a^2+B_{\rm rot}j_b^2
&=\frac{A_{\rm rot}+B_{\rm rot}}{4}
(j_+j_-+j_-j_+)\\
&\quad+\frac{A_{\rm rot}-B_{\rm rot}}{4}
(j_+^2+j_-^2).
\end{align*}

再用

\[
\frac12(j_+j_-+j_-j_+)=\mathbf j^2-j_c^2,
\]

可将一般不对称陀螺写成

\begin{equation}
H_{\rm rot}
=\frac{A_{\rm rot}+B_{\rm rot}}{2}\mathbf j^2
+\left(C_{\rm rot}-\frac{A_{\rm rot}+B_{\rm rot}}2\right)j_c^2
+\frac{A_{\rm rot}-B_{\rm rot}}4(j_+^2+j_-^2).
\label{eq:asymmetric-top-k-mixing}
\end{equation}

在 \(|JKM\rangle\) 基中，对角矩阵元为

\begin{equation}
\langle JKM|H_{\rm rot}|JKM\rangle
=\frac{A_{\rm rot}+B_{\rm rot}}2J(J+1)
+\left(C_{\rm rot}-\frac{A_{\rm rot}+B_{\rm rot}}2\right)K^2.
\label{eq:asymmetric-top-diagonal-element}
\end{equation}

而

\[
j_+|JKM\rangle
=\sqrt{J(J+1)-K(K+1)}\,|J,K+1,M\rangle,
\]

连续作用两次得到

\begin{align*}
j_+^2|JKM\rangle
&=C^+_{JK}|J,K+2,M\rangle,\\
C^+_{JK}
&=\sqrt{[J(J+1)-K(K+1)]
[J(J+1)-(K+1)(K+2)]}.
\end{align*}

因此唯一的非零带外矩阵元是

\begin{equation}
\langle J,K+2,M|H_{\rm rot}|JKM\rangle
=\frac{A_{\rm rot}-B_{\rm rot}}4 C^+_{JK},
\label{eq:asymmetric-top-offdiagonal-element}
\end{equation}

以及其厄米共轭 \(K\to K-2\)。所以非对称性不是任意混合所有 \(K\)，而是把固定 \(J,M\) 子空间分裂成若干只通过 \(\Delta K=\pm2\) 相连的块。\(K\) 不再是好量子数，但 \(K\) 的奇偶性等离散对称信息仍然约束矩阵结构。

常用 Ray asymmetry 参数记为

\begin{equation}
\kappa_{\rm R}
=\frac{2B_{\rm rot}-A_{\rm rot}-C_{\rm rot}}
{A_{\rm rot}-C_{\rm rot}},
\qquad -1\le\kappa_{\rm R}\le1.
\label{eq:ray-asymmetry-parameter}
\end{equation}

其中 \(\kappa_{\rm R}=-1\) 是长球极限，\(\kappa_{\rm R}=+1\) 是扁球极限。一般 \(\kappa_{\rm R}\) 下没有单个闭式 \(E_{JK}\)，应在每个固定 \(J,M\) 的有限矩阵中对角化 \eqnrefs{eq:asymmetric-top-k-mixing}。

同一个转子同时允许两种群作用。空间固定转动改变分子在实验室中的整体取向，可以看成从左侧作用于旋转矩阵；体固定转动则改变分子坐标架的内部标记，对应右侧作用。设 \(L_i\) 与 \(R_i\) 分别是 \(SO(3)\) 上的左不变和右不变向量场。Lie 群上的一般结果是

\begin{equation}
[L_i,L_j]=\varepsilon_{ijk}L_k,
\qquad
[R_i,R_j]=-\varepsilon_{ijk}R_k,
\qquad
[L_i,R_j]=0.
\label{eq:left-right-invariant-algebra}
\end{equation}

负号来自右不变场的群交换子顺序与 Lie 代数元素乘法顺序相反。量子角动量由 \(-\ii\hbar\) 乘相应微分生成元得到，因此

\begin{equation}
[J_X,J_Y]=+\ii\hbar J_Z,
\qquad
[J_a,J_b]=-\ii\hbar J_c,
\label{eq:body-fixed-anomalous-commutator}
\end{equation}

以及循环置换。

\begin{proof}[解]
取 Lie 代数元素 \(X,Y\)。Baker--Campbell--Hausdorff 公式给出左作用的无穷小群回路
\[
 \ee^{\epsilon X}\ee^{\eta Y}
 \ee^{-\epsilon X}\ee^{-\eta Y}
 =\exp\!\left(\epsilon\eta[X,Y]
 +O(\epsilon^2\eta,\epsilon\eta^2)\right).
\]
因此 \([L_X,L_Y]=L_{[X,Y]}\)。对右作用，同样的微分算符组合使群元按反序累积，对应回路为
\[
 \ee^{-\eta Y}\ee^{-\epsilon X}\ee^{\eta Y}\ee^{\epsilon X}
 =\exp\!\left(-\epsilon\eta[X,Y]+O(\epsilon^2\eta,\epsilon\eta^2)\right).
\]
因此 \([R_X,R_Y]=-R_{[X,Y]}\)。

量子角动量分别定义为
\(J_i=-\ii\hbar L_i\) 与 \(J_a=-\ii\hbar R_a\)。于是
\begin{equation*}
 [J_X,J_Y]=+\ii\hbar J_Z,\qquad
 [J_a,J_b]=-\ii\hbar J_c ,
\end{equation*}
即\eqnrefs{eq:body-fixed-anomalous-commutator}。
\end{proof}

这个负号只记录体固定坐标架采用右作用，并不表示物理旋转代数失效。实验室系与分子本体系生成元彼此对易，所以 \(M\) 和 \(K\) 可同时作为对称陀螺的量子数。不对称项含 \(j_+^2+j_-^2\)，只连接 \(K\) 与 \(K\pm2\)；因此 \(K\) 本身不再守恒，而 \(K\) 的奇偶性仍然守恒。

\section{不可约极化率张量与排列选择定则}
对远离真实电子共振、且电场频率足够高到可以对光学周期平均的非极性分子，主要相互作用来自诱导偶极

\[
\mathbf p(t)=\boldsymbol\alpha\mathbf E(t),
\]

其中 \(\boldsymbol\alpha\) 是分子极化率二阶张量。若实电场写成 \(\mathbf E(t)=\mathbf E_0\cos\omega t\)，周期平均相互作用为

\begin{equation}
\overline V_{\rm int}
=-\frac14\mathbf E_0^T\boldsymbol\alpha\mathbf E_0.
\label{eq:polarizability-cycle-averaged}
\end{equation}

对线性分子

\[
\boldsymbol\alpha
=\alpha_\perp\mathbf 1
+\Delta\alpha\,\hat{\mathbf n}\hat{\mathbf n},
\qquad
\Delta\alpha=\alpha_\parallel-\alpha_\perp.
\]

若线偏振沿实验室 \(Z\) 轴，\(\theta\) 是分子轴和 \(Z\) 的夹角，则

\begin{equation}
V(\theta)
=-\frac14E_0^2
[\alpha_\perp+\Delta\alpha\cos^2\theta].
\label{eq:linear-molecule-alignment-potential}
\end{equation}

利用

\[
\cos^2\theta
=\frac13+\frac23P_2(\cos\theta),
\]

可把 \eqnrefs{eq:linear-molecule-alignment-potential} 分成

\begin{align*}
V(\theta)
&=-\frac14E_0^2
\left(\alpha_\perp+\frac{\Delta\alpha}{3}\right)
-\frac{E_0^2\Delta\alpha}{6}P_2(\cos\theta).
\end{align*}

第一项是秩-0 标量，只给所有转动态共同的 AC Stark 位移；真正改变转动态组成的是第二个秩-2 项。这就是为什么取向的角动量选择定则由二阶球张量而不是一阶偶极张量决定。

一般地，两个秩-1 场张量的耦合写成

\[
[E^{(1)*}\otimes E^{(1)}]^{(K)}_Q,
\qquad K=0,1,2.
\]

对非磁、无手性且极化率对称的普通分子，反对称秩-1 部分消失，保留 \(K=0,2\)。两个同阶不可约张量的旋转标量积定义为

\[
T^{(K)}\!\cdot U^{(K)}
=\sum_{Q=-K}^{K}(-1)^Q T^{(K)}_QU^{(K)}_{-Q},
\]

于是相互作用可以统一写成

\begin{equation}
H_{\rm int}
=-\frac12\sum_{KQ}(-1)^Q
[E^{(1)*}\otimes E^{(1)}]^{(K)}_Q
\alpha^{(K)}_{-Q},
\qquad K=0,2,
\label{eq:polarizability-tensor-coupling}
\end{equation}

其中整体数值因子随复振幅定义改变，但张量结构与选择定则不变。

对线性刚性转子，场沿 \(Z\) 时只有 \(P_2(\cos\theta)\) 需要计算。利用

\[
P_2(\cos\theta)=\sqrt{\frac{4\pi}{5}}Y_{20}(\theta,\phi)
\]

和三个球谐函数的 Gaunt 积分，可以得到

\begin{equation}
\begin{aligned}
&\langle J'M'|P_2(\cos\theta)|JM\rangle\\
&\quad=(-1)^{M'}
\sqrt{(2J'+1)(2J+1)}
\begin{pmatrix}
J'&2&J\\0&0&0
\end{pmatrix}
\begin{pmatrix}
J'&2&J\\-M'&0&M
\end{pmatrix}.
\end{aligned}
\label{eq:alignment-p2-matrix-element}
\end{equation}

第二个 3-j 符号要求

\[
M'=M,
\]

三角条件要求 \(|J-J'|\le2\le J+J'\)，而第一个 3-j 符号还要求 \(J+J'+2\) 为偶数。因此线偏振取向的主要选择定则是

\begin{equation}
\Delta M=0,
\qquad
\Delta J=0,\pm2.
\label{eq:alignment-selection-rule}
\end{equation}

这与电偶极跃迁常见的 \(\Delta J=\pm1\) 根本不同，因为这里驱动系统的是秩-2 极化率各向异性。

\section{绝热驱动、脉冲激发与转动回生}

同一个势 \eqnrefs{eq:linear-molecule-alignment-potential} 在不同脉冲时间尺度下产生两种很不相同的物理。对线性转子，场自由能谱为

\[
E_J=B_eJ(J+1),
\]

因此典型转动时间尺度为 \(\hbar/B_e\)，完整量子回生时间为

\begin{equation}
T_{\rm rev}=\frac{\pi\hbar}{B_e}.
\label{eq:rotational-revival-time}
\end{equation}

若脉冲包络变化时间 \(\tau_p\) 满足

\[
\tau_p\gg\frac{\hbar}{|E_n-E_m|}
\]

并且更严格的绝热条件

\begin{equation}
\epsilon_{nm}^{\rm ad}
\equiv
\frac{\hbar|\langle n(t)|\dot H(t)|m(t)\rangle|}
{|E_n(t)-E_m(t)|^2}
\ll1
\label{eq:rotational-adiabatic-parameter}
\end{equation}

对所有重要耦合态成立，分子会跟随瞬时 pendular eigenstate，称绝热取向。若场也绝热关闭，系统大致回到原来的无场转动态，强取向不会长期保留。

相反，当

\[
\tau_p\ll\frac{\hbar}{B_e},
\]

脉冲期间自由转动来不及显著演化，可以把激光看成一个瞬时“冲击”。忽略脉冲内的 \(H_{\rm rot}\) 后，演化算符近似为

\begin{equation}
U_{\rm kick}
\simeq
\exp\!\left[
\ii P\cos^2\theta
\right],
\qquad
P=\frac{\Delta\alpha}{4\hbar}
\int_{-\infty}^{\infty}E_0^2(t)\,\dd t.
\label{eq:alignment-kick-strength}
\end{equation}

\(P\) 是脉冲激发的无量纲强度。脉冲结束后，多个 \(J\) 态按照 \(E_J\) 自由积累相位；正是这些由 \eqnrefs{eq:alignment-selection-rule} 逐级建立的 \(\Delta J=\pm2\) 相干在 \(T_{\rm rev}\) 的整数或分数倍附近重新相位匹配，产生无场取向回生。

\begin{proof}[解]
脉冲期间的精确传播子为
\[
U(t_+,t_-)
=\mathcal T\exp\left[-\frac{\ii}{\hbar}
\int_{t_-}^{t_+}\bigl(H_{\rm rot}+V(t)\bigr)\,\dd t\right].
\]
若脉冲只显著占据一组有限的转动态，并且该子空间内满足
\[
\epsilon_{\rm imp}
\equiv
\frac{\tau_p\,\Delta E_{\rm rot}^{(\rm acc)}}{\hbar}\ll1,
\]
则在脉冲宽度内自由转动造成的相对相位可以作为受控小量。Magnus 展开的首项给出
\[
U(t_+,t_-)
=\exp\left[-\frac{\ii}{\hbar}\int V(t)\,\dd t\right]
+O(\epsilon_{\rm imp}),
\]
而 \(J\)-无关的标量 AC Stark 项只贡献整体相位。代入\eqnrefs{eq:linear-molecule-alignment-potential}便得到
\[
U_{\rm kick}
=\exp\!\left(\ii P\cos^2\theta\right)
+O(\epsilon_{\rm imp}).
\]

脉冲后线性刚性转子自由演化。对任意整数 \(J\)，
\[
U_0(T_{\rm rev})|JM\rangle
=\exp\!\left[-\frac{\ii B_eJ(J+1)}{\hbar}
\frac{\pi\hbar}{B_e}\right]|JM\rangle.
\]
由于相邻整数的乘积 \(J(J+1)\) 必为偶数，
\[
\exp[-\ii\pi J(J+1)]=1,
\]
所以所有被脉冲激发的 \(J\) 分量在
\[
T_{\rm rev}=\frac{\pi\hbar}{B_e}
\]
处同时恢复相位，得到完整回生。分数回生则来自同一二次谱相位在 \(T_{\rm rev}/q\) 处按同余类重新组织，而不是额外引入的新动力学机制。
\end{proof}

\begin{note}[Alignment versus orientation]
\eqnrefs{eq:linear-molecule-alignment-potential} 只含 \(\cos^2\theta\)，在 \(\hat{\mathbf n}\to-\hat{\mathbf n}\) 下不变，所以它只能使分子轴取向而不能区分"头"和"尾"。真正区分分子头尾的定向需要奇宇称耦合，例如永久偶极与静电场的 \(-\boldsymbol\mu\cdot\mathbf E\) 或合适双色场产生的奇阶有效相互作用。
\end{note}

\section{转动光谱实验与分子结构测定}\label{sec:rotation-experiments}

本节给出转动光谱的典型实验参数、不对称陀螺的数值实例，以及激光取向的实验参数。

纯转动光谱（微波区）由 \(\Delta J=1\) 电偶极跃迁给出，频率
\[
\tilde\nu(J\to J+1)=2\tilde B_e(J+1)-4\tilde D_e(J+1)^3
\]
永久偶极矩 \(\mu\ne0\) 是纯转动光谱活性的必要条件：同核双原子（H\(_2\)、N\(_2\)、O\(_2\)）无纯转动红外/微波谱。

\begin{center}
\begin{tabular}{lcccc}
\hline
分子 & \(\tilde B_e\) (cm\(^{-1}\)) & \(\mu\) (D) & \(T_{\rm rev}=\pi/(2\tilde B_ec)\) (ps) & 跃迁频率 \(J=0\to1\) (GHz) \\
\hline
HCl & 10.593 & 1.08 & 0.79 & 634.8 \\
CO & 1.931 & 0.112 & 4.33 & 115.8 \\
NO & 1.702 & 0.153 & 4.91 & 102.2 \\
OCS & 0.203 & 0.715 & 41.2 & 12.2 \\
HCN & 1.478 & 3.00 & 5.65 & 88.7 \\
\hline
\end{tabular}
\end{center}

\begin{exbox}[CO 的纯转动光谱]
CO 的 \(J=0\to1\) 跃迁在 115.8\,GHz（\(\lambda\approx2.59\,\mathrm{mm}\)），是天体化学中探测星际 CO 分子的标志性谱线。前几条线：\(J=1\to2\) 在 231.6\,GHz，\(J=2\to3\) 在 347.4\,GHz——等间距 \(2\tilde B_ec\approx115.8\,\mathrm{GHz}\) 是刚性转子的特征。离心畸变使间距随 \(J\) 略减：\(J=10\to11\) 的间距为 115.5\,GHz，偏差 \(\sim0.3\%\)。
\end{exbox}

水分子的主轴转动常数 \(A_{\rm rot}/(hc)=27.877\,\mathrm{cm}^{-1}\)，\(B_{\rm rot}/(hc)=14.512\,\mathrm{cm}^{-1}\)，\(C_{\rm rot}/(hc)=9.285\,\mathrm{cm}^{-1}\)，Ray 参数
\[
\kappa_{\rm R}=\frac{2B_{\rm rot}-A_{\rm rot}-C_{\rm rot}}{A_{\rm rot}-C_{\rm rot}}
=\frac{29.024-27.877-9.285}{27.877-9.285}\approx-0.437
\]
中等不对称，\(K\) 混合显著。

\begin{exbox}[\(J=1\) 水分子能级]
不对称陀螺 \(J=1\) 有 3 个能级（\(M\) 简并不计）。取 \(K\) 绕 c 轴（最小转动常数 \(C_{\rm rot}\) 的方向）投影，\(K=0\) 为偶块、\(K=\pm1\) 为奇块。\(K\) 偶块：\(\langle J_c^2\rangle=0\)、\(\langle J_a^2+J_b^2\rangle=J(J+1)=2\)，故 \(E_{1,0}=A_{\rm rot}+B_{\rm rot}=42.389\,\mathrm{cm}^{-1}\)。\(K=\pm1\) 奇块是 \(2\times2\) 矩阵：对角元 \(d=\frac{A_{\rm rot}+B_{\rm rot}}{2}+C_{\rm rot}=30.480\,\mathrm{cm}^{-1}\)，非对角元 \(\Delta=\frac{A_{\rm rot}-B_{\rm rot}}{2}=6.683\,\mathrm{cm}^{-1}\)——来自 \(J_a^2-J_b^2=\frac12(J_+^2+J_-^2)\) 把 \(K=+1\) 与 \(K=-1\) 耦起来（\(\langle 1,-1|J_a^2-J_b^2|1,1\rangle=1\)）。本征值为 \(d\pm\Delta=37.162\) 和 \(23.797\,\mathrm{cm}^{-1}\)。三个 \(J=1\) 刚性转子能级为 23.797、37.162、42.389\,\(\mathrm{cm}^{-1}\)；实际分子另有很小的离心畸变修正。
\end{exbox}

\eqnrefs{eq:alignment-kick-strength} 的脉冲强度 \(P\) 取决于极化率各向异性 \(\Delta\alpha\) 和脉冲能量。典型参数：

\begin{center}
\begin{tabular}{lcccc}
\hline
分子 & \(\Delta\alpha\) (\AA\(^3\)) & \(B_e/h\) (GHz) & \(T_{\rm rev}\) (ps) & 典型 \(P\) \\
\hline
N\(_2\) & 0.69 & 59.9 & 8.4 & 5--10 \\
O\(_2\) & 0.60 & 43.4 & 11.5 & 4--8 \\
CO\(_2\) & 1.85 & 11.7 & 42.8 & 10--20 \\
I\(_2\) & 6.0 & 1.12 & 448 & 50--100 \\
\hline
\end{tabular}
\end{center}

\begin{exbox}[\(I_2\) 的转动回生]
I\(_2\) 的 \(T_{\rm rev}\approx448\,\mathrm{ps}\) 是实验上最容易观测的回生之一。Stapelfeldt 实验用 800\,nm、100\,fs 脉冲（\(\tau_p\ll\hbar/B_e\approx140\,\mathrm{ps}\)）激发，脉冲强度 \(P\approx60\)。在 \(T_{\rm rev}/4\approx112\,\mathrm{ps}\) 观测到 4 束波包（四分之一回生），在 \(T_{\rm rev}/2\approx224\,\mathrm{ps}\) 观测到 2 束（半回生），在 \(T_{\rm rev}\approx448\,\mathrm{ps}\) 完全复原。分数回生的存在直接验证了 \(E_J\propto J(J+1)\) 的二次依赖——若能级是线性的，回生不会出现分数结构。
\end{exbox}

\begin{proof}[解]
\eqnrefs{eq:asymmetric-top-k-mixing} 只含对角项以及
\(K\to K\pm2\) 的矩阵元，所以 \(K\) 的奇偶性守恒。再利用
哈密顿量在 \(K\leftrightarrow-K\) 下不变，对 \(K>0\) 定义
\begin{equation}
|J,K,\sigma\rangle=\frac{1}{\sqrt2}\bigl(|J,K\rangle+\sigma\,|J,-K\rangle\bigr),
\qquad
K>0,\ \sigma=\pm1,
\label{eq:wang-transform}
\end{equation}
取 \(K=0\) 态归入 \(\sigma=+1\) 子空间。定义 Wang 宇称算符 \(\mathsf P|J,K\rangle=|J,-K\rangle\)，则 \(|J,K,\sigma\rangle\) 正是 \(\mathsf P\) 的本征态（本征值 \(\sigma\)）。哈密顿量在 \(K\leftrightarrow-K\) 下不变意味着 \([H,\mathsf P]=0\)，于是
\[
 \langle J,K,\sigma|H|J,K',\sigma'\rangle
 =\sigma\sigma'\langle J,K,\sigma|H|J,K',\sigma'\rangle,
\]
即 \((1-\sigma\sigma')\langle\cdots\rangle=0\)：除非 \(\sigma=\sigma'\)，矩阵元严格为零（不只是正比），故
\[
 \langle J,K,\sigma|H|J,K',\sigma'\rangle\propto\delta_{\sigma\sigma'}.
\]
而 \(K-K'\) 必为偶数。因此固定 \(J,M\) 后，
\begin{equation}
H=H_{E^+}\oplus H_{E^-}\oplus H_{O^+}\oplus H_{O^-},
\label{eq:wang-block-diagonalization}
\end{equation}
式 \eqnrefs{eq:wang-block-diagonalization} 中，\(E/O\) 表示偶/奇 \(K\)，上标表示 Wang 宇称。该分块来自精确离散对称性，不依赖不对称性大小。
\end{proof}

数值对角化时只需分别处理这四个子块。这个结论比低 \(J\) 的逐个显式解更有用：它同时给出能级的离散对称标签，并阻止不同 Wang 对称类之间产生避免交叉。
